\documentclass[11pt,a4paper]{article}
\input{../../sablony/pripravy/priprava.sty}
\usepackage{tikz}
\nastavPripravu{1. ročník}{Souhrnné opakování}{131}{Algebraický servis}
\begin{document}
\blok{Úvod a připomenutí}
Opakování algebraických vzorců, rozkladů, úprav mnohočlenů a podmínek lomených výrazů. Motiv: \emph{kontrola vadného výpočtu} — studenti mají chybu pojmenovat a opravit. V každém lomeném výrazu napiš podmínky \textbf{před} krácením.
\blok{Společně řešené příklady}
\begin{prikladReseniBox}[title={Příklad 1}]
Upravte $(2x-3)^2-(x+1)(x-1)$.
\tcblower
\[
\begin{aligned}
(2x-3)^2-(x+1)(x-1)&=4x^2-12x+9-(x^2-1)\\
&=3x^2-12x+10
\end{aligned}
\]

\end{prikladReseniBox}
\begin{prikladReseniBox}[title={Příklad 2}]
Rozložte na součin $6x^2-x-2$.
\tcblower
\[
\begin{aligned}
6x^2-x-2&=(2x+1)(3x-2)
\end{aligned}
\]

\end{prikladReseniBox}
\begin{prikladReseniBox}[title={Příklad 3}]
Upravte $\dfrac{x^2-5x+6}{x^2-9}$ a určete všechna nepřípustná $x$.
\tcblower
\[
\begin{aligned}
x^2-9&=(x-3)(x+3)\ne0\\
x&\ne3,\quad x\ne-3\\
\frac{x^2-5x+6}{x^2-9}&=\frac{(x-2)(x-3)}{(x-3)(x+3)}\\
&=\frac{x-2}{x+3}\quad(x\ne\pm3)
\end{aligned}
\]

\end{prikladReseniBox}
\begin{prikladReseniBox}[title={Příklad 4}]
Žák napsal $(x-3)^2=x^2-9$. Proč to není správně?
\tcblower
\[
\begin{aligned}
(x-3)^2&=(x-3)(x-3)\\
&=x^2-6x+9\\
&\ne x^2-9\quad\text{obecně}
\end{aligned}
\]

\end{prikladReseniBox}
\blok{Samostatné procvičování}
\begin{zadaniReseniBox}[title={Samostatná práce 1}]
Upravte $(x+4)^2-(x-4)^2$.
\tcblower
\[
\begin{aligned}
(x+4)^2-(x-4)^2&=16x
\end{aligned}
\]

\end{zadaniReseniBox}
\begin{zadaniReseniBox}[title={Samostatná práce 2}]
Rozložte $4x^2-25$ na součin.
\tcblower
\[
\begin{aligned}
4x^2-25&=(2x-5)(2x+5)
\end{aligned}
\]

\end{zadaniReseniBox}
\begin{zadaniReseniBox}[title={Samostatná práce 3}]
Zjednodušte $\dfrac{x^2-1}{x-1}$.
\tcblower
\[
\begin{aligned}
x&\ne1\\
\frac{x^2-1}{x-1}&=\frac{(x-1)(x+1)}{x-1}=x+1
\end{aligned}
\]

\end{zadaniReseniBox}
\blok{Volitelná rezerva a výzva}
\poznamka{Výzva: vysvětli, proč se čitatelem a jmenovatelem nesmí „zkrátit“ sčítané členy. Např. $(x+2)/(x+3)$ není $2/3$.}
\blok{Domácí úkol}
\begin{zadaniReseniBox}[title={Domácí úkol}]
\begin{enumerate}[label=\alph*),itemsep=2pt]
\item Rozviňte $(3x-2)^2$.
\item Upravte $2(x-1)^2-(x+1)(x-3)$.
\item Rozložte $4x^2-12x+9$.
\item Rozložte $2x^2+5x+2$.
\item Zjednodušte $\dfrac{x^2-4}{x^2+3x+2}$ včetně podmínek.
\item Určete všechna $k\in\mathbb R$, pro něž je $x^2+kx+16$ druhou mocninou lineárního dvojčlenu s reálnými koeficienty.
\end{enumerate}
\tcblower
\begin{enumerate}[label=\alph*),itemsep=2pt]
\item
\[
\begin{aligned}
(3x-2)^2&=9x^2-12x+4
\end{aligned}
\]
\item
\[
\begin{aligned}
2(x^2-2x+1)&-(x^2-2x-3)\\
&=x^2-2x+5
\end{aligned}
\]
\item
\[
\begin{aligned}
4x^2-12x+9&=(2x-3)^2
\end{aligned}
\]
\item
\[
\begin{aligned}
2x^2+5x+2&=(2x+1)(x+2)
\end{aligned}
\]
\item
\[
\begin{aligned}
x^2+3x+2&=(x+1)(x+2)\ne0\\
x&\ne-1,-2\\
\frac{x^2-4}{x^2+3x+2}&=\frac{x-2}{x+1}
\end{aligned}
\]
\item
\[
\begin{aligned}
x^2+kx+16&=(x\pm4)^2\\
&=x^2\pm8x+16\\
K&=\{-8;8\}
\end{aligned}
\]
\end{enumerate}
\end{zadaniReseniBox}
\poznamka{Čas: 5 + 17 + 15 + 8 minut. Důsledně odlišuj \emph{rovnost výrazů na přípustném oboru} od rovnosti pro všechna reálná čísla.}
\end{document}
